\section{Result Comments}

The aerodynamic coefficient set shows that drag increases with Mach number and angle of attack, as expected for a finned slender body. The pitching moment changes significantly with alpha, and the reference condition at Mach 0.9 gives \(C_{M\alpha}<0\), indicating a restoring pitching trend around the selected moment reference.

In the alpha versus delta result, the \(2^\circ\) deflection produces a limited alpha response, with a peak near \(4.7^\circ\). The approximate peak ratio is therefore \(\alpha_{peak}/\delta \approx 2.35\). Physically, the control surface generates enough moment to perturb the attitude, while the response remains inside the tabulated aerodynamic range. Beta remains smaller than alpha because the command mainly excites the longitudinal plane; the nonzero beta is attributed to the coupled 6DOF equations and remains bounded during the short test.

The missile mass decreases during the booster burn, and the CG moves as propellant is consumed. In the short 2.8 s test the CG moves slightly aft, from 2.833 m to 2.938 m, because the consumed booster propellant is ahead of the initial overall CG. For complete burnout, the estimated CG would move forward to 1.856 m after the sustainer propellant is also consumed. This CG variation changes the moment arms and therefore the dynamic response, which is why the common moment reference and the time-varying CG table must be kept consistent.
